\documentclass[11pt]{article}
\usepackage[T1]{fontenc}
\usepackage[utf8]{inputenc}
\usepackage{lmodern,amsmath,amssymb,amsthm,mathtools,longtable}
\usepackage[margin=1in]{geometry}
\usepackage[hidelinks]{hyperref}
\hypersetup{pdftitle={Supplement to Certified whole-domain bounds and equality barriers for the odd theta–Weil form},
 pdfauthor={Leonard van Hemert},
 pdfsubject={Version 1.0.0; mathematical preprint and certified proof scope}}
\usepackage{microtype}
\setlength{\emergencystretch}{1em}
\makeatletter
\renewcommand*\l@section{\@dottedtocline{1}{0em}{2.8em}}
\renewcommand*\l@subsection{\@dottedtocline{2}{2.8em}{3.6em}}
\makeatother
\newtheorem{theorem}{Theorem}[section]
\newtheorem{lemma}[theorem]{Lemma}
\newtheorem{proposition}[theorem]{Proposition}
\newtheorem{corollary}[theorem]{Corollary}
\newcommand{\Ran}{\operatorname{ran}}
\newcommand{\Dom}{\operatorname{Dom}}
\newcommand{\dd}{\,d}
\renewcommand{\thesection}{S\arabic{section}}
\renewcommand{\theequation}{\thesection.\arabic{equation}}
\numberwithin{equation}{section}
\title{Supplement to \emph{Certified whole-domain bounds and equality barriers
for the odd theta--Weil form}}
\author{Leonard van Hemert}
\date{13 October 2026 --- Version 1.0.0}
\begin{document}\maketitle
\begin{abstract}
This supplement supplies the analytic domain and tail arguments, finite
matrix error formulas, and certificate correspondence for the main article.
It also records the secondary signed-action and coupling results. All
norms are in the original complex Hilbert space unless a coefficient
norm is explicitly specified. Exact stored binary numbers define trial
maps; floating solves and spectral decompositions select those numbers
but do not establish inverses, rank, or positivity. The finite arithmetic
certificates and their source files accompany this document.
\end{abstract}
\tableofcontents
\setcounter{section}{-1}
\section{Foundations, domains, and primitive tails}\label{foundations}
We use the definitions $\Phi,w,p,q,G,\mathcal V,A_W,E_8,J,J_8$ of
Section~1 of the main article. Write $W=w^2$, $\Omega=\widehat W$,
$\Xi=\widehat\Phi$, $Q=\|q\|^2$, $\tau=\|p\|^2$ and
$S=I-q\otimes q/Q$. Pairings are sesquilinear, and $f\otimes f$
denotes the positive map $u\mapsto\langle u,f\rangle f$, with the
corresponding consistent convention for adjoints. The whole-line
normalization equals $2\dd x$ on the positive half-line for products
of odd functions. In particular the positive-frequency measure is
$r_8(t)\dd t/\pi$, not $r_8(t)\dd t/(2\pi)$.

\subsection{Closed realization and domain-preserving kernels}
The original signed realization and unconditional pole-null jet
identity are \cite[Lemma 5.1]{Theta}; its closed form and compact
core are proved in \cite[Lemma 2.2]{Theta}. We spell
out the additional facts needed for this article. The positive symbol has
the absolutely convergent representation
\[
 b(t)=\sum_{j\ge0}\frac{(t/2)^2}
 {(j+1/4)((j+1/4)^2+(t/2)^2)}.
\]
Multiplication by bounded $w$ followed by the closed multiplier
$b(D)^{1/2}$ is closed: convergence of its input and output implies
convergence of $wu$ and its multiplier image, so closedness of the
multiplier applies. Its squared norm defines the closed form $g$.
The prime majorant below and the finite-rank pole show that the
remaining signed form is a bounded selfadjoint perturbation. Thus
$A_W$ is selfadjoint on $\Dom(G)$ and has form domain $\mathcal V$.

Writing $a_j=2j+1/2$ gives
\[
 b'(t)=4t\sum_{j\ge0}\frac{a_j}{(a_j^2+t^2)^2},\qquad t>0.
\]
The step-two unimodal sum is bounded by half its integral plus twice
its supremum. Consequently $b'(t)\le9/t$ for $t\ge1$, while
$b(1)\le\tfrac14\sum(j+1/4)^{-3}<17$. Hence
\begin{equation}\label{logbound}
 0\le b(t)\le17+9\log(1+|t|).
\end{equation}
For each fixed frequency, $W\sin(tx)$, $W\sinh(x/2)$ and every
fixed derivative of $\Phi$ are Schwartz. The entire theta series and
its differentiated series converge with a common superexponential
majorant on the real line. Equation~\eqref{logbound} therefore places
the corresponding physical columns in $\Dom(G)$.

The score $q=\Phi'/w$ and first pole-null jet
$h=(\Phi'''-\Phi'/4)/w$ also belong to $\Dom(G)$ and to
$\ker A_W$. Here is the full domain argument underlying that statement.
For every fixed $j$, termwise differentiation gives
$|\Phi^{(j)}(x)|\le D_j e^{2j|x|}\Phi(x)$ with finite $D_j$:
a polynomial in $y=\pi n^2e^{2|x|}$ multiplies each theta atom, and
the residual sum is bounded by
$\sum n^{4+2j}e^{-\pi(n^2-1)}<\infty$. Division by the positive
first theta atom bounds the ratios and shows that $q,h$ are square
integrable with all required theta-weighted derivatives. Their
weighted transforms are
\[
 \widehat{wq}(z)=-iz\Xi(z),\qquad
 \widehat{wh}(z)=iz(z^2+1/4)\Xi(z).
\]
They vanish at every actual nontrivial zero, without assuming its
ordinate real. The second also vanishes at the poles $z=\pm i/2$.
Apply the signed explicit formula first to compactly cut off physical
tests and then remove the cutoffs. The relative-theta bounds just given
control every differentiated cutoff error; contour shifts within
$|\Im x|<\pi/4$ and the zero-count bound below dominate the paired
zero sum. The full geometric side contains the contact, every prime
power and both poles before odd restriction. Thus its mixed pairing
with either vector vanishes on the compact form core. By continuity
in the form norm, it vanishes on $\mathcal V$. The representation
theorem gives $A_Wq=A_Wh=0$ and their operator-domain membership.
For the score, the nonzero pole values are included in this signed
identity; they are not discarded as if the score were pole-null.

More generally, let $K$ be a closed subspace of $\ker A_W$. If
$A_W=G-D$ with $D$ bounded, then $GP_K=DP_K$ is bounded. It follows
that $P_K$ preserves both $\Dom(G)$ and $\mathcal V$ and
\begin{equation}\label{kerneltransfer}
 \mathfrak a(u)=\mathfrak a((I-P_K)u),\qquad u\in\mathcal V.
\end{equation}
Indeed the spectral projection onto the kernel reduces $A_W$;
any subprojection inside that zero block also reduces it. Alternatively
expand the form and use its vanishing mixed pairings. Since
$|d_{8,+}-b|\le |c_\Gamma|+\|k_8\|_\infty$, the operators $E_8$
and $G$ differ by a bounded selfadjoint operator. They have identical
operator and form domains. The positive compressed operator
$\mathcal E$ on $q^\perp$ has domain $\Dom(G)\cap q^\perp$ and
value $SE_8u$: the missing pairing with $q$ is bounded because
$E_8q\in H$. This justifies every resolvent and projection used in
the main article.

\subsection{Complete prime tail}
Put $q_0=e^{-\pi}$, $d=1+16q_0^3/(1-6q_0)$ and $a=3/(2\pi)$.
The inequalities $j^4\le16\,6^{j-2}$ and $j^2-1\ge j+1$ for $j\ge2$
retain every theta atom and give
\[
 \Phi(x)\le4\pi^2e^{9|x|/2-\pi e^{2|x|}}(d-ae^{-2|x|}).
\]
For a shift $\log n$, write $A=e^{2|x|}$, $B=e^{2|x-\log n|}$,
$t=\sqrt{AB}\ge n$. Using $A+B\ge2t$ and
$(d-a/A)(d-a/B)\le(d-a/t)^2$, and dropping only the positive
cosh denominator, gives
\[
 w(x)w(x-\log n)\le4\pi^2t^{9/4}e^{-\pi t}(d-a/t).
\]
This decreases for $t\ge2$, since its logarithmic derivative is at
most $9/8+1/6-\pi<0$. Therefore
\begin{equation}\label{primebound}
 \|J-J_8\|\le8\pi^2\sum_{n\ge9}
       \Lambda(n)n^{7/4}e^{-\pi n}(d-a/n)
       <t_8:=\frac{32237}{10^{12}}.
\end{equation}
For an exact rational verification take
$\pi_+=355/113$, $q_+=2160875733/(5\cdot10^{10})$ and
$\delta_+=348730539/(2\cdot10^{11})$. Then $\pi<\pi_+$,
$e^{-\pi}<q_+$ and $d\le1+\delta_+$. Since
$\Lambda(n)n^{7/4}\le n^3$, the last sum is at most
\[
 8\pi_+^2(1+\delta_+)\,
       \frac{9^3q_+^9}{1-q_+(10/9)^3}<\frac{32237}{10^{12}}.
\]
The rational inequalities follow by cross multiplication; the elementary
$\pi$ and exponential endpoints follow from Machin's formula with
alternating remainders and the exponential Taylor remainder. The same
summable majorant proves operator-norm convergence before the split.
Thus this is a whole operator bound, not a bound on a sampled prime sum.

\subsection{Full gamma cutoff and complete cells}
The exact principal-log integral for $\psi$ \cite[5.9.13]{DLMF}
yields, for $|t|\ge2$,
\[
 m_\Gamma(t)\ge\log(|t|/(2\pi))-211/300.
\]
To see the explicit error, $\coth s\le1/s+s/3$ follows by comparing
Taylor coefficients in $(3+s^2)\sinh s-3s\cosh s$. Hence
$1/(1-e^{-x})-1/x\le1/2+x/12$. The integral at $5/4+it/2$ bounds
$|\psi(z)-\log z|$ by $34/75$; recurrence back to $1/4+it/2$
costs at most $1/t^2\le1/4$. The six actual prime powers have
$\|k_8\|_\infty<13/2$. Since $\log(10000/(2\pi))>29/4$,
\begin{equation}\label{cutoff}
 d_8(t)>7/150\quad(|t|\ge10000),\qquad r_8(t)=0\quad(|t|\ge10000).
\end{equation}
For reproducible exact checks, the square-root lower endpoints for
$n=2,3,4,5,7,8$ are respectively $7/5,5/3,2,11/5,13/5,14/5$.
Logarithm upper endpoints may be checked by positive exponential
Taylor sums. The logarithm inequality uses $e<11/4$,
$e^{1/4}\le4/3$ and $\pi<22/7$.

On $[0,10000]$, all $80000$ closed cells of width $1/8$ are covered.
Outward evaluation classifies $79524$ as zero and $476$ as possibly
positive. The latter are refined to width $1/256$; their $11476$
possibly nonzero subcells lie in $2905$ cells of width $1/64$.
For a subcell with midpoint $c$, halfwidth $h$ and lower endpoint
$a>0$, the full gamma derivative bound is
\[
 L(a)=\min(10,2/a^2+\pi/(2a)),\qquad L(0)=10.
\]
It follows by differentiating the entire digamma series and comparing
the remaining positive tail with its integral. Thus the center
$\Re\psi(1/4+ic/2)-\log\pi$ plus $[-Lh,Lh]$ encloses the full
symbol. Evaluating all six cosines on the whole interval and taking
the positive part gives upper and lower rectangles for $r_8$.
Every upper mass is rounded upward to a dyadic endpoint of denominator
$2^{80}$, then to the exact upward binary64 endpoint used to define $B$.
Their sum is below $14$. Zero cells, refined cells and \eqref{cutoff}
cover the entire frequency line by evenness.

\subsection{Theta and zero envelopes}
For $x\ge0$, the full lattice gives
$\Phi(x)\le80e^{9x/2-3e^{2x}}$ and
$W(x)\le160e^{4x-3e^{2x}}$. Contour shifting gives
\begin{equation}\label{fouriertail}
 |\Omega(t)|,|v(t)|\le5e^{-|t|/4},\qquad
 |t\Xi(t)|\le5|t|e^{-|t|/4},
\end{equation}
where $v(t)=2\int_0^\infty W(x)\sinh(x/2)\sin(tx)\dd x$.
For completeness, on $|\Im x|\le1/4$ use
$\pi\cos(1/2)>21/8$ and
$|\cosh((x+i/4)/2)|\ge(127/128)\cosh(x/2)$.
With $q=1/13$ and
\[
 S_0\le q+q^4/(1-q^3),\quad
 S_2\le q+4q^4/(1-3q^3),\quad
 S_4\le q+16q^4/(1-6q^3),
\]
the complete strip bounds are
$|W(z)|\le(256/127)(40S_4+19S_2)<10$ and
\[
 \int|\Phi(x+i/4)|\dd x,\quad
 \int|W(x+i/4)\sinh((x+i/4)/2)|\dd x
       \le(128/127)(15S_2+25S_0)<5.
\]
The analogous $W$ integral is smaller. They follow by
$t=e^{2|x|}$ and integration of the positive exponential-polynomial
majorants. Vertical contour ends vanish superexponentially, proving
\eqref{fouriertail}. These estimates also give the looser weighted
$1/8$-strip integrals bounded by $120$ used below.

The unconditional zero count, with both signs and multiplicity, is
bounded for $L\ge e^2$ by $2L\log L$. The explicit bound of \cite[Corollary 1.2]{HSW} is
$|N_+(T)-\frac{T}{2\pi}\log(T/(2\pi e))|
\le0.1038\log T+0.2573\log\log T+9.3675$ for $T\ge e$.
Twice this upper bound is at most $2T\log T$ for $T\ge e^2$:
use $\log\log T\le\log T$, bound the main coefficient by $1/\pi$,
and check the endpoint; the difference thereafter increases.
This retains multiplicity and both signs. Stieltjes integration then gives the complete tail
\begin{equation}\label{zerotail}
 Z_a(L)=e^{-aL}\left[2L\log L+
 \frac{2(\log L+1)}a+\frac4{a^2L}\right],\qquad a>0.
\end{equation}
Indeed integrate the counting majorant against $ae^{-at}$ on
$[L,\infty)$, twice integrate by parts, and use
$\int_L^\infty e^{-at}/t\dd t\le e^{-aL}/(aL)$.
The bound includes possible off-line zeros: if
$\rho=1/2+i\gamma$, then $|\Im\gamma|<1/2$. Cauchy--Schwarz and
the complete theta envelope give
$|\widehat{wu}(\gamma)|\le11\|u\|$. Thus exponentially decaying
samples may be paired against every Hilbert test in the full signed
zero formula, rather than only other trial vectors.

\section{Continuum features and the midpoint comparison}\label{midpoint}
Put $f_t=S[w\sin(tx)]$ and $p_Y=Sp$. On the common coefficient
space $L^2(r_8(t)\dd t/\pi)\oplus\mathbb C$, the continuum
feature map has positive product
\[
 R_q=SR_8S=\int_0^{10000} f_t\otimes f_t\,r_8(t)\dd t/\pi
             +2p_Y\otimes p_Y.
\]
Let $I_i=[k_i/64,(k_i+1)/64]$, $t_i=(2k_i+1)/128$ and let
$\mu_i^+$ be the exact stored upward masses. Define
$Be_i=\sqrt{\mu_i^+}f_{t_i}$, $Be_p=\sqrt2p_Y$ and $R_m=BB^*$.
The square roots are mathematical square roots of exact dyadics.

Because $|x|\le2|\sinh(x/2)|$ and $\tau<3/500$,
\[
 \|f_t-f_{t_i}\|^2\le |t-t_i|^2\int x^2W(x)\dd x
 \le4\tau|t-t_i|^2.
\]
The full mass is below $14$ and $|t-t_i|\le1/128$, whence
\begin{equation}\label{delta}
 \int\|f_t-f_{t_i}\|^2\,r_8(t)\dd t/\pi
       \le\delta_0:=21/1024000.
\end{equation}
For any $a>0$, vector Young's inequality on every whole cell proves
\begin{equation}\label{youngmid}
 R_q\le(1+a)R_m+(1+1/a)\delta_0I.
\end{equation}
Replacing the true cell mass by its upper mass is legitimate only
on the nonnegative midpoint term. The pole is exact in both maps;
enlarging it by $1+a$ is conservative.

Let $G_0(i,j)=[\Omega(t_i-t_j)-\Omega(t_i+t_j)]/2$,
$G_0(p,i)=v(t_i)$, $G_0(p,p)=\tau$,
$b_i=-t_i\Xi(t_i)$ and $b_p=-1/4$.
For $D_\mu=\operatorname{diag}(\sqrt{\mu_i^+},\sqrt2)$,
\[
 B^*B=D_\mu(G_0-bb^*/Q)D_\mu\le\Gamma,
\]
where $\Gamma$ replaces $Q$ by $12$ and the pole norm by $3/500$.
The difference is positive semidefinite, since $Q<12$ and
$\tau<3/500$. All difference, sum, score and pole entries remain.

\section{Dense capture and arithmetic spectral enclosures}\label{capture}
The exact dyadic $2906\times56$ matrix $Z$ and positive diagonal
$\Lambda$ satisfy the certified full dense residual bound
\[
 \|\Gamma-Z\Lambda Z^*\|<3/12500.
\]
Let $P$ be the actual coefficient projection onto $\Ran Z$. Since
$(I-P)Z=0$ and $\Gamma\ge0$,
$0\le B(I-P)B^*\le(3/12500)I$. Once $N=Z^*Z>0$ is established,
$P=ZN^{-1}Z^*$. Put $C=BZ$ and $R_{56}=(11/10)CN^{-1}C^*$.
Using \eqref{youngmid} with $a=1/10$ and the complete prime tail gives
\begin{equation}\label{rest}
 \mathfrak a(z)\ge E_8(z)-\langle R_{56}z,z\rangle-\rho\|z\|^2,
 \quad \rho=t_8+11\delta_0+\frac{11}{10}\frac3{12500}
       =\frac{979236349}{2000000000000}.
\end{equation}
This holds throughout $\mathcal V\cap q^\perp$. No physical mixed
block is removed by the coefficient-space decomposition.

The center entry error for $\Gamma$ is below
$3.124405\cdot10^{-9}$ in operator norm. Near Fourier arguments use
outward dyadic endpoint pairs; far arguments use \eqref{fouriertail}.
Each short scalar entry expression receives $10^{-12}$ for endpoint
conversion, products, square-root scales and symmetrization. Absolute
row sums receive the positive summation guard below. The full center
Frobenius residual is $0.0002391007522$ to the displayed precision;
its guarded upper endpoint is below $0.000239106<3/12500$.
The entire dense residual is evaluated, not just retained columns.

Here are the common arithmetic rules used throughout the certificate.
For binary64 unit roundoff $u=2^{-53}$, write
$\gamma_n=nu/(1-nu)$. A length-$m$ matrix dot product receives
$\gamma_{2m+10}$ times its absolute-product sum. Positive sums used
to bound those absolute products are themselves enlarged by their
corresponding $\gamma$ factor. A Frobenius norm receives a serial
square/sum bound for all entries, followed by an outward square root.
Symmetrization and short final scalar formulas receive their stated
additional absolute allowance. Nonnegative row sums are divided by
$1-\gamma_n$ and rounded upward. Every intermediate is finite and
away from overflow; outward nextafter steps cover conversion underflow.
No unbounded cancellation is hidden in a relative sum error.

For a real symmetric center $A$, let the exact dyadic locator $V,D$
have orthogonality defect $o=\|V^*V-I\|<1$, full residual
$r=\|AV-VD\|$ and $L=\|D\|$. Its exact polar unitary
$U=V(V^*V)^{-1/2}$ obeys
\begin{equation}\label{polar}
 \|U^*AU-D\|\le\frac r{\sqrt{1-o}}
 +2L\sqrt{1+o}\left(\frac1{\sqrt{1-o}}-1\right).
\end{equation}
Expand $AU-UD$ into $(AV-VD)(V^*V)^{-1/2}$ and the remaining
commutator. Functional calculus bounds the latter by twice $L$
times the distance of $(V^*V)^{-1/2}$ from $I$. Add any whole
center error $\|A_{\rm true}-A\|$ before using the spectral endpoint.
This proves that eigensolvers supply locators only. The later small
comparisons use exact-input 192-bit Arb products and Frobenius bounds
in \eqref{polar}.

\section{Physical Gram and genuine rank}\label{physical}
The coefficient and physical metrics are different:
\[
 N=Z^*Z,\qquad H_C=C^*C=Z^*B^*BZ.
\]
The full coefficient defect is below $3.627\cdot10^{-11}$ after
dot and row guards. Hence
\begin{equation}\label{ngram}
 (1-10^{-8})I\le N\le(1+10^{-8})I.
\end{equation}
The physical Gram uses the actual $Q$ and $\tau$, not their loose
upper substitutes. The complete norm integrals are
\[
 Q=2\int_0^\infty\frac{\Phi'^2\cosh(x/2)}{\Phi}\dd x,
 \qquad \tau=2\int_0^\infty\frac{\Phi\sinh^2(x/2)}{\cosh(x/2)}\dd x.
\]
Put $t=e^{2x}$ and $y_n=\pi n^2t$. The first two atom polynomials
are $4y^2-6y$ and $-8y^3+30y^2-15y$. For $y\ge3$, positivity of
the first and the complete lattice sum imply
$|\Phi'|\le80t\Phi$. Consequently both doubled integral tails
beyond $x=2$ are bounded by
\[
 512000\int_{54}^\infty t^4e^{-3t}\dd t<2^{-180}.
\]
On $[0,2]$ the atoms $n\ge9$ satisfy
$\sum_{n\ge9}n^re^{-\pi t(n^2-1)}\le2\,9^re^{-240}$ for $r=4,6$.
Relative to the retained positive eight-atom sum, the omitted function
and derivative are each below $2^{-300}$. Expanding the quotient
shows that their total integral effect is below $2^{-250}$.
The common allowance is therefore $2^{-250}+2^{-180}$.

Native meromorphic Arb integration on $[0,2]$ encloses the finite
integrals, with absolute and relative tolerance $2^{-100}$. A possible
zero in a denominator ball produces a nonfinite ball rather than an
analyticity assertion. The finite returned balls, enlarged by both
complete allowances, give
\begin{align*}
 Q&=11.0676657471428749463191678229\ \pm6.97\cdot10^{-29},\\
 \tau&=0.00571290852978025013036375087\ \pm4.55\cdot10^{-30}.
\end{align*}
The actual whole Gram error is below $3.138\cdot10^{-9}$. Compression
adds less than $5.785\cdot10^{-11}$, and multiplying the input error
by $\|Z\|^2\le1+10^{-8}$ gives total below $3.2\cdot10^{-9}$.
Score products of either sign use both $Q$ endpoints. The exact pole
entry is $2(\tau-1/(16Q))$ and the cross term is
$v(t_i)+b_i/(4Q)$ before mass scaling.

For the saved physical spectral locator, safe rational upper errors
are $d=32/10^{10}$, $o=43/10^{12}$ and $r=115/10^{13}$ in
\eqref{polar}. They prove
\begin{equation}\label{physicalbounds}
 H_C\ge\frac{21}{10^8}N,
 \qquad H_C<\frac{32}{125}I.
\end{equation}
It follows that the actual physical rank is $56$ and the coefficient
preimage cost is at most $\kappa_B\le\sqrt{10^8/21}<2182.178903$.
No enlarged enriched-trial inverse or rank is inferred from this.

\section{Fourier grids and the complete low-zero frame}\label{zeros}
\subsection{Outward quadrature}
The basic grid evaluates $\Omega,\Omega'$ at $j/64$ through $100$
and $v,-t\Xi(t)$ at the half-step frequencies. Native 192-bit Arb
Legendre roots and weights of order $256$ are mapped to $[0,3]$.
The Bernstein ellipse of parameter $r=23/20$ has imaginary part less
than $1/4$, real part below $16/5$ and $|z|<4$. The strip estimates
in S0 give the common integrand bound $M=240\cdot3^{25}$.
Chebyshev coefficients satisfy $|a_j|\le2Mr^{-j}$; positivity and
polynomial exactness through degree $511$ give the doubled integral
error
\[
 E_{256}=24M r^{-512}/(1-r^{-1})+32\,2^{-128}+2^{-1160}.
\]
The second term covers all omitted theta atoms; the last covers the
whole physical tail. For example $t=e^{2x}$ reduces that tail to
$160e^{-1200}(729^2/3+2\cdot729/9+2/27)<2^{-1175}$.
Every transform ball is enlarged before conversion to outward
integer endpoints over $2^{80}$. For the score multiply the enclosed
$\Xi(t)$ by the exact frequency and retain its enlarged radius.

Arbitrary zero arguments use an order-$20$ Taylor grid at $j/8$
through $128$, with native quadrature order $320$. Derivative $j$
costs $4^j$ in the ellipse bound; its complete base error is
\[
 E_{320}=24\cdot240\cdot3^{32}r^{-640}/(1-r^{-1})
                     +32\,2^{-128}+2^{-1100}.
\]
For $j\le19$, the physical derivative tails are bounded by
$160\int_{400}^\infty t^6e^{-3t}\dd t<2^{-1100}$. At displacement
$|h|\le1/16$ the entire Taylor remainder is at most
$120(2|h|)^{20}$, using the exponential moment bound from S0.
Horner products, coefficient radii, frequency addition and native zero
location errors are priced. Global derivative bounds are
$|\Omega'|<11/100$ and $|v'|<3/250$. Far arguments retain
\eqref{fouriertail} with a downward dyadic argument reserve; both
sum and difference arguments are included through $6000$.

\subsection{Actual zeros, signs and multiplicity}
Native FLINT 3.6.0 \texttt{zeta\_zeros} and \texttt{zeta\_nzeros}
\cite{FLINT} certify $2469$ disjoint actual positive roots on the
critical line and the full multiplicity count $N(3000)=2469$.
Their last upper ordinate is below $3000$; hence the roots exhaust
the count and are simple within this finite interval. The intervals
are stored outward over $2^{128}$. Both signs are retained by oddness.
Define $V_0e_j=w\sin(\gamma_jx)$, $A_0=V_0^*C$ and $G_z=V_0^*V_0$.
At actual zeros the score transform vanishes exactly. Thus
\[
 S_{3000}=2V_0V_0^*,\quad L=C^*S_{3000}C=2A_0^*A_0,
 \quad O=C^*S_{3000}^2C=4A_0^*G_zA_0.
\]
The complete sample entries use
$\sqrt{\mu_i^+}[\Omega(\gamma-t_i)-\Omega(\gamma+t_i)]/2$
and pole entry $\sqrt2v(\gamma)$. All $2469^2$ zero correlations and
all $2469\cdot2906$ zero/feature correlations are included.
The full certified errors are
\[
 \|\Delta A_0\|_F\le2\cdot10^{-12},\quad
 \|\Delta G_z\|\le3.2\cdot10^{-11},\quad
 \|\Delta L\|\le10^{-11},\quad\|\Delta O\|\le3.1\cdot10^{-10}.
\]
To certify a small positive gap, the rectangular factor $A_0$ is
checked directly. Its singular-factor locator has left and right
orthogonality defects at most $10^{-11},2\cdot10^{-13}$ and full
residual at most $3\cdot10^{-14}$. Subtracting this residual and
the complete sample error from
$\sigma_{\min}\sqrt{1-o_L}\sqrt{1-o_R}$ gives at least
$3.2586\cdot10^{-8}$. Consequently $L\ge10^{-15}I$.
This argument does not subtract the larger additive $L$-error from
an unverified tiny eigenvalue.

\subsection{The omitted zeros as a whole action bound}
The actual largest midpoint is $T=313317/128$. If $f=B\xi$ with
minimal coefficient preimage, $\|\xi\|\le\kappa_B\|f\|$.
The sine sum on the $1/8$ strip has size at most
$4e^{T/8}\|\xi\|$, and the pole coefficient at most $2\|\xi\|$.
Using the full strip integrals bounded by $120$ yields, at every
actual zero,
\[
 |\widehat{wf}(\gamma)|\le1000\kappa_B
       e^{(T-|\Re\gamma|)/8}\|f\|.
\]
The subtracted score vanishes at these zeros. Squaring this bound and
using \eqref{zerotail} bounds the form remainder; pairing it with the
whole Hilbert evaluation bound $11\|v\|$ bounds the operator remainder:
\[
 e_L=10^6\kappa_B^2e^{T/4}Z_{1/4}(L),\qquad
 d_L=11000\kappa_Be^{T/8}Z_{1/8}(L).
\]
At $L=3000$ these satisfy $e_L<3\cdot10^{-43}$ and
$d_L<1.4\cdot10^{-18}$. All multiplicities, both signs and possible
high off-line zeros remain. The mixed formula extends from compact
core tests by the same theta domination as S0, proving
$\|(A_W-S_{3000})|_{\Ran C}\|<d_L$.
Together with \eqref{physicalbounds} and $L\ge10^{-15}I$ this gives
$\mathfrak a|_{\Ran C}\ge(38/10^{16})I$. It supplies no positive
lower bound for the infinite complement.

\section{True masses and the weighted deficit}\label{mass}
For each subcell, the lower rectangle from S0 gives an outward
lower mass $\ell_i$ on the cell $I_i$, rounded downward over $2^{80}$.
Thus $0\le\ell_i\le\mu_i:=\int_{I_i}r_8(t)\dd t/\pi\le\mu_i^+$.
The exact deficits $d_i=\mu_i^+-\ell_i$ have sum below $0.29826603$.
For the true-mass midpoint operator $R_t$, including the identical pole,
\[
 0\le R_m-R_t\le R_d:=\sum_i d_i f_{t_i}\otimes f_{t_i}.
\]
A trace estimate is valid but unnecessarily large. Instead every
entry of the actual score Gram is bounded by the positive upper Gram
$\overline G=G_0-bb^*/12$. Its full entry intervals use the near
Fourier endpoints and all far envelopes. If $A_{ij}\ge|\overline G_{ij}|$
and $w_i>0$, the weighted Schur estimate is
\[
 \|R_d\|\le\max_i\frac{\sum_j A_{ij}d_jw_j}{w_i}.
\]
The weights are exact positive integer multiples of $2^{-30}$, bounded
below by $1/1024$. Their method of selection imposes no convergence
premise. Entry errors receive $10^{-12}$, and the positive row products,
sums and division receive $\gamma_{2\cdot2905+20}$, an outward step
and a final exact $10^{-15}$ allowance. The guarded maximum is below
$0.008552658376$, giving
\begin{equation}\label{deficit}
 \|R_m-R_t\|<427633/50000000=0.00855266.
\end{equation}
The pole deficit is zero; this removes no physical pairing with a pole
test vector because the remaining feature map is kept in full.

\section{Whole-space midpoint error and frame action}\label{bridge}
Regard the continuum and true-mass midpoint feature maps as maps
on the same $L^2$ input space. By \eqref{delta} their difference has
norm at most $\sqrt{\delta_0}$. Expanding their positive products gives
\[
 \|R_q-R_t\|\le2\sqrt{\|R_t\|\delta_0}+\delta_0.
\]
The dense capture and exact factor provide
$\|R_t\|\le\|R_m\|\le(1+10^{-8})\max\Lambda+3/12500
<0.255985513$. Thus the whole-space bound is
\begin{equation}\label{epsilonr}
 \|R_q-R_m\|\le\epsilon_R:=
 2\sqrt{0.255985513\delta_0}+\delta_0+0.00855266.
\end{equation}
On the score quotient, the exact domain identity is
$\mathcal E=A_W+R_q+S(J-J_8)S$. Therefore on $\Ran C$,
\[
 \|\mathcal E-(S_{3000}+R_m)\|\le\epsilon_R+t_8+d_{3000}
 <0.013155646076930=:\eta.
\]
This last action bound is restricted to $\Ran C$; it is never applied
to newly enriched zero directions. The whole-space bound
\eqref{epsilonr}, by contrast, is independent of a trial space.

\section{Complete frame residual and center errors}\label{residual}
This section supplies the finite realization of main Lemma~2.1.
For $s>0$, let $Y_0=CX$ with exact dyadic $X$. Put
$G_B=B^*B$, $Q_b=V_0^*B$ and $G_z=V_0^*V_0$.
Then, with $E_m=S_{3000}+R_m$,
\[
 r=C-(E_m+s)Y_0=BD_1+V_0D_2,
 \quad D_1=Z(I-sX)-G_BZX,\quad D_2=-2A_0X,
\]
\begin{align*}
 r^*r={}&D_1^*G_BD_1+D_1^*Q_b^*D_2
               +D_2^*Q_bD_1+D_2^*G_zD_2.
\end{align*}
Every term is a full physical Gram. In particular $R_mC$ is not
replaced by $CH_C$, and the residual is not projected into $\Ran C$.
The full $Q_b$ Frobenius error is below $7.215\cdot10^{-13}$;
all other input errors are those of S3--S4.

If binary coefficient maps differ from the exact $D_1,D_2$ by
$e_1,e_2$, their physical map error is at most
$e_d=\sqrt{\|G_B\|}e_1+\sqrt{\|G_z\|}e_2$.
Let $S_c$ approximate the Gram of the binary map with error $r_e$.
Then the exact residual Gram has error bounded above by
\[
 s_e=r_e+2\sqrt{\lambda_{\max}(S_c)+r_e}\,e_d+e_d^2.
\]
This prices input maps before the matrix Young comparison.
If $h_e,l_e,m_e$ bound respectively the physical Gram, zero compression
and midpoint compression errors, the remaining exact-center errors obey
\begin{align*}
 q_e&\le h_e\|X\|^2+
       \|Q_c-X^*H_cX\|_F,\\
 k_e&\le2h_e\|X\|+(l_e+m_e+sh_e)\|X\|^2\\
 &\hspace{12mm}+\|K_c-(H_cX+X^*H_c-X^*(L_c+M_c+sH_c)X)\|_F.
\end{align*}
The displayed center differences are computed using exact binary
inputs and exact rational $s$ in Arb. Thus shift conversion, all inner
and outer products and symmetrization are included. Each raw numerical
guard is enlarged by $10^{-12}$; a further $10^{-6}$ is added to the
final comparison price.

The complete upper prices, with $s=\mu-\rho$, are
\[
 \begin{array}{c|c|c}
 \mu&\text{Young parameter}&\text{complete price}\\\hline
 1/20&2/3&100157725018/125000000000\\
 1/25&1/2&450253581783/500000000000
 \end{array}
\]
They are strictly below the respective targets
\[
 \frac{400631}{500000}(1-10^{-8}),\qquad
 \frac{225127}{250000}(1-10^{-8}),
\]
with exact margins $9592169/(5\cdot10^{13})$ and
$20685723/(25\cdot10^{12})$.
Factorization in main Lemma~2.1 and \eqref{rest} give the stated
frame-only comparisons. Their role is to exhibit the effect of keeping
matrix correlation and to supply the residual directions used next.

\section{The equality barrier and kernel-aware projection}\label{barriersupp}
This supplies the scalar certificate for main Theorem~5.1.
With $h=(\Phi'''-\Phi'/4)/w$, set
$a=\langle h,q\rangle/Q$, $k=h-aq$ and $D=\|k\|^2$.
S0 proves $k\in\Dom(G)\cap\ker A_W\cap q^\perp$.
Its transform $it(t^2+1/4+a)\Xi(t)$ is not identically zero, so $D>0$.
The degree-five atom polynomial for $wh$ is
\[
 -32y^5+360y^4-1056y^3+810y^2-90y.
\]
For $y\ge3$ its absolute value is at most $301y^5$. The full lattice
comparison
$\sum n^{10}q^{n^2-1}\le1+1024q^3/(1-58q^5)<2$, $q<1/20$,
gives $|wh|/\Phi<20000e^{6x}$.
Both doubled tails beyond $x=2$ of the norm and score pairing are
bounded by
\[
 10^{11}\int_{54}^\infty t^8e^{-3t}\dd t<2^{-140}.
\]
On $[0,2]$, the omitted degree-ten lattice satisfies
$\sum_{n\ge9}n^{10}e^{-\pi t(n^2-1)}\le2\,9^{10}e^{-240}$.
The three relative errors for $\Phi,\Phi',wh$ are below $2^{-270}$.
Expanding the two quotient integrands bounds their integral errors
by $12\,2^{85}2^{-270}<2^{-170}$. The common allowance $2^{-130}$
is added to each meromorphic Arb integral before projection.
This gives
\[
 \langle h,q\rangle=73.9389206536588932761816\ \pm6.26\cdot10^{-23},
\]
\[
 \|h\|^2=54152.6907182449014086321392\ \pm3.28\cdot10^{-23}.
\]
In particular $a>6$ and $D<55000$.

Every term of $R_m$ is nonnegative. The $635$ stored cells with
$t_i\le20$ give the valid lower sum
\[
 \langle R_mk,k\rangle\ge
 \sum_{t_i\le20}\mu_i^+t_i^2(t_i^2+25/4)^2\Xi(t_i)^2>1875.
\]
The exact dyadic masses and native complex $\Xi$ evaluations yield
\[1876.7393730639281947402372447352786\]
with radius below $2\cdot10^{-31}$. Thus $r=\langle R_mk,k\rangle/D>3/88$.
The common feature-map triangle inequality gives
$\langle R_qk,k\rangle\le
(\sqrt{\langle R_mk,k\rangle}+\sqrt{\delta_0}\|k\|)^2$.
Using $\mathfrak a(k)=0$, the dense capture and full prime tail proves
\begin{align*}
 \frac{\langle R_{56}k,k\rangle-E_8(k)}D&\ge f(r),\\
 f(r)&:=r/10-2\sqrt{\delta_0r}-\delta_0\\
     &\hspace{12mm}-(11/10)(3/12500)-t_8.
\end{align*}
The derivative is positive for $r\ge3/100$, and exact outward
arithmetic gives $f(3/88)>7/5000$. This is the strict main barrier.
Testing $R_{56}\le\theta(\mathcal E+s)$ on $k$ gives
$\theta s>7/5000+(1-\theta)E_8(k)/D$ for $0<\theta\le1$.

For an enlarged kernel removing $k$, let $\Pi=I-k\otimes k/D$,
$B'=\Pi B$ and $C'=\Pi C$. Exact projection gives
\[
 B'^*B'=B^*B-(B^*k)(B^*k)^*/D,\quad
 C'^*C'=C^*C-(C^*k)(C^*k)^*/D.
\]
Since $V_0^*k=0$, zero sampling and zero/feature cross Grams remain
unchanged. Injectivity of $V_0^*C$ implies the projected rank is still
$56$. Domain preservation follows from S0. An action comparison
obtained by subtracting $A_WC$ before projection retains a term
$d\|CX\|$, as well as the new continuum error on $C'X$; it cannot
replace $\|CX\|$ by $\|C'X\|$ without a separate estimate. These
identities establish a permissible changed certificate, not a new
numerical improvement or a barrier on that larger quotient.

\section{Admissible enrichment and a constructive action estimate}\label{enrich}
Let $Y=Bb+V_0v:\mathbb C^{56}\to H$ use the supplied exact binary
coefficient arrays. Each column is in $\Dom(G)\cap q^\perp$ by S0.
A solve may choose the arrays but imposes no accuracy, rank or enlarged
Gram-inverse premise. In particular the old frame-only error $\eta$
is not available on $Y$.
Put $\ell_8=m_\Gamma-k_8$ and $L_8=M_w\ell_8(D)M_w$.
The pointwise identity $d_{8,+}=\ell_8+r_8$ and S6 give, on the
whole operator domain,
\begin{equation}\label{fullbridge}
 \mathcal E=SL_8S+R_m-2p_Y\otimes p_Y+\Delta_R,
 \qquad\Delta_R=R_q-R_m,\quad\|\Delta_R\|\le\epsilon_R.
\end{equation}
This replaces the unpriced high-zero action on the enriched columns
by a smooth multiplier action.

For one column with base representation
$y_0=w\sum_a c_a\sin(ax)+c_pw\sinh(x/2)$, write
$y=Sy_0$ and $\alpha=\langle y_0,q\rangle/Q$. Then
\begin{equation}\label{exactfourier}
 \widehat{wy}(t)=i\left\{
 \sum_a c_a\frac{\Omega(t-a)-\Omega(t+a)}2
             +c_pv(t)+\alpha t\Xi(t)\right\}.
\end{equation}
The coefficients include the mathematical square roots of the actual
masses. For $V_0$ columns the score overlaps vanish at the actual
zero ordinates. The pole overlap is exactly $-1/4$.

Here is an explicit finite action enclosure. Put $a_0=1/4$ and
$b_0=1/8$. For $|s|\le1/4$, the complete theta series proves
\begin{equation}\label{strip20}
 |\Omega(t+is)|,\ |\widehat{W\sinh(x/2)}(t+is)|,
 |\Xi(t+is)|\le20e^{-|t|/8}.
\end{equation}
For an explicit bound use $\pi\cos(1/4)>5/2$ and
$|\cosh((x+i\theta)/2)|\ge\tfrac12\cosh(x/2)$.
On $x\ge0$, the weighted $W$ majorant is
\[
 \sum_n[160n^4e^{17x/4}+76n^2e^{9x/4}]
                         e^{-(5/2)n^2e^{2x}}.
\]
Substitute $t=e^{2x}$ and bound its fractional powers by $t^2,t$.
Their integrals are at most $106/125$ and $14/25$, times
$e^{-5n^2/2}$. The lattice sum is at most $25/246$, so both
half-lines cost
$2[80(106/125)+38(14/25)](25/246)<20$.
The other two functions have smaller bounds. Contour shifting by
$1/8$ gives \eqref{strip20}, including both physical tails.

The analytic continuation
$m_\Gamma(z)=[\psi(1/4+iz/2)+\psi(1/4-iz/2)]/2-\log\pi$
has no poles in this strip. The complete trigamma series bounds its
vertical derivative by $36$, and the six-term cosine bound together
with \eqref{logbound} gives
\[
 |\ell_8(t+is)|\le200+10|t|,\qquad |s|\le1/4.
\]
Regarding $c_a,c_p,\alpha$ as coefficient rows for the whole map $Y$,
a common real and strip-line $L^1$ bound for
$\ell_8(z)\widehat{wY}(z)$ is
\begin{equation}\label{envelopeA}
 A=20\left[\sum_a(4480+160|a|)\|c_a\|_2
                +4480\|c_p\|_2+47200\|\alpha\|_2\right].
\end{equation}
The constants come from the complete moments
$\int e^{-|t|/8}=16$, $\int|t|e^{-|t|/8}=128$ and
$\int t^2e^{-|t|/8}=2048$.

Let $h>0$, $T=Kh$, $P=2\pi/h$ and physical radius $R>0$.
If $L_k$ denotes the exact row
$\ell_8(kh)\widehat{wY}(kh)$ and $\widehat L_k$ its exact binary
approximant, define
\[
 \widehat g(x)=\frac h{2\pi}\sum_{|k|\le K}\widehat L_ke^{-ikhx},
 \quad\widehat J=S[w\widehat g],
\]
\[
 t_h=\frac h{2\pi}\sum_{|k|>K}\|L_k\|_2,\quad
 n_h=\frac h{2\pi}\sum_{|k|\le K}\|L_k-\widehat L_k\|_2,
 \quad B_h=\frac h{2\pi}\sum_{k\in\mathbb Z}\|L_k\|_2.
\]
Every infinite sum is bounded by shifted exponential envelopes in
\eqref{strip20}, or the sharper real bounds \eqref{fouriertail}.
Geometric series and their first two derivatives sum their polynomial
factors. No primitive frequency is replaced by a smaller presumed band.
Poisson summation and contour inversion give the complete map bound
\begin{align}\label{actionbound}
 \|SL_8Y-\widehat J\|\le{}&\|w\|_2
 \left[\frac A\pi e^{a_0R}
 \frac{e^{-a_0P}}{1-e^{-a_0P}}+t_h+n_h\right]\\
 &+\left[\frac A{2\pi}+B_h\right]
                      \|w1_{|x|>R}\|_2.\nonumber
\end{align}
Indeed the inverse transform has row norm at most
$Ae^{-a_0|x|}/(2\pi)$. Its periodic aliases on $[-R,R]$ sum to
the first displayed term. Outside that interval use $A/(2\pi)$
and $B_h$ on the two full inverse transforms, then multiply by $w$.
Frequency omission and node error cost respectively $t_h,n_h$.
The projection $S$ is contractive. This proves a whole output-space
map estimate, including components outside every trial span.

\section{The executed action certificate and its full Grams}\label{action}
The parameters are $h=1/20$, $T=3300$, $R=3$. Near Fourier arguments
use the complete Taylor enclosures in S4, with Horner, coefficient,
zero-location and mass-square-root errors. The signed multiplier is
evaluated at every exact grid point by outward digamma and six cosine
balls. Score values through $128$ use order-$320$ quadrature with the
complete base error in S4; beyond $128$ they use \eqref{fouriertail}.
The following upward endpoints include all nodes and tails:
\begin{center}
\begin{tabular}{lr}
Quantity & Upper bound\\\hline
Strip row integral $A$ & $1.001591\cdot10^9$\\
Alias before $\|w\|_2$ & $1.532843\cdot10^{-5}$\\
Finite node/location/rounding error before $\|w\|_2$
 & $1.797275\cdot10^{-8}$\\
Omitted frequencies before $\|w\|_2$ & $10^{-88}$\\
Both physical tails after $w$ & $10^{-157}$\\
Finite midpoint-action map error & $3.366659\cdot10^{-6}$\\
Complete map error $g$ & $1.421821\cdot10^{-5}$\\
Trial norm $\|Y\|$ & $2.159290$
\end{tabular}
\end{center}
Here $\|w\|_2\le1/\sqrt2$ and
$\|w1_{|x|>3}\|_2\le2^{-550}$. The finite midpoint action is
$Bd$, where $d$ is the exact binary approximation to
$G_Bb+Q_b^*v$ with its pole row removed. Its coefficient error is
priced as a physical map using $\|B\|$ and the full input Gram errors.
Set $\widehat A=\widehat J+Bd$. Then
\begin{equation}\label{actionpremise}
 \|\mathcal EY-\widehat A-\Delta_RY\|\le g.
\end{equation}

The full physical action Gram is essential. For the smooth part it is
\[
 (h/(2\pi))^2\widehat L^*
       [\Omega((k-l)h)]\widehat L
 - (h/(2\pi))^2(b_0^*\widehat L)^*(b_0^*\widehat L)/Q.
\]
The Toeplitz product uses native outward complex DFT and inverse at
length $262144$, on exact binary input. Every returned radius is
retained. Kernel values through $128$ use the complete Taylor
certificate; the omitted kernel has its whole geometric $L^1$ price.
The pole, score, $Bd$ and all physical/Fourier cross products receive
their full point and absolute-product guards.
For
$G_Y=Y^*Y$, $CY=C^*Y$, $YA=Y^*\widehat A$,
$CA=C^*\widehat A$, $AG=\widehat A^*\widehat A$, the resulting
operator uncertainties have upward endpoints
\begin{center}
\begin{tabular}{lr}
$G_Y$ & $0.014115739$\\
$CY$ & $4.979513\cdot10^{-5}$\\
$YA$ & $5.596759\cdot10^{-6}$\\
$CA$ & $1.982337\cdot10^{-8}$\\
$AG$ & $2.263431\cdot10^{-9}$
\end{tabular}
\end{center}
These are matrix uncertainties; they are not silently treated as a
whole-map error. In particular $AG$ includes the entire outside output.

\section{Shifted residual theorem and the finite tradeoffs}\label{tradeoffs}
Put $T_s=\mathcal E+sI$, $s=\mu-\rho>0$. Completing the square gives
\[
 C^*T_s^{-1}C=K+R^*T_s^{-1}R\le K+R^*R/s,
 \quad R=C-T_sY,\quad K=C^*Y+Y^*C-Y^*T_sY.
\]
This uses $T_s\ge sI$ and the operator-domain membership of all trial
columns, without an unshifted gap. Define
\begin{align*}
 K_s&=CY+CY^*-sG_Y-(YA+YA^*)/2,\\
 S_s&=H_C-CA-CA^*-s(CY+CY^*)+AG+s(YA+YA^*)+s^2G_Y.
\end{align*}
Write the exact error as $\Delta_RY+D$, $\|D\|\le g$.
Its quadratic contribution is at most
$\epsilon_RG_Y+g\|Y\|I$, and for $\tau>0$ its squared Gram is
at most $(1+\tau)\epsilon_R^2G_Y+(1+1/\tau)g^2I$.
A second Young inequality separating this error from the whole
residual $C-\widehat A-sY$ gives, for every $\beta>0$,
\begin{align}\label{enrichedcriterion}
 C^*T_s^{-1}C\le{}&K_s+
 \left[\epsilon_R+\frac{(1+1/\beta)(1+\tau)\epsilon_R^2}s\right]G_Y
 +\frac{1+\beta}sS_s\\
 &+\left[g\|Y\|_{\rm up}
       +\frac{(1+1/\beta)(1+1/\tau)g^2}s\right]I.\nonumber
\end{align}
The chosen $\tau=g/(\epsilon_R\|Y\|_{\rm up})$ is positive.
Each primitive matrix uncertainty from S10 is multiplied by its
actual absolute scalar coefficient in the assembly of $K_s,S_s$.
Exact shift conversion and all scalar products receive their stated
rounding guards. The finite spectral comparison uses \eqref{polar}
with exact-input 192-bit Arb residuals. Multiplication of the entire
right side by $11/10$ and division of its upper endpoint by
$n_0=1-10^{-8}$ yield an upward rational $\theta$ on the $10^{-12}$
lattice. By \eqref{ngram} and the factorization norm identity,
\[
 R_{56}\le\theta T_s,\qquad
 \mathfrak a(u)\ge(1-\theta)E_8(z)-[\rho+\theta s]\|z\|^2,
 \quad z=(I-P_K)u.
\]
The final transfer is exactly \eqref{kerneltransfer}; it does not
recompute a numerical certificate on a changed frame.

There are $231$ tested rational parameter pairs with individually
enclosed endpoints. Some have $\theta\ge1$ and do not retain a
strictly positive energy coefficient. Among those with $\theta<1$,
the least scalar loss is the first row below. The stated theorem uses
exact terminating coefficients; scalar losses in the table are rounded
upward for presentation. The exact first loss is given below the table.
\begin{center}
\begin{tabular}{c|c|c|c}
$\mu$ & $\beta$ & $1-\theta$ & scalar loss upper\\\hline
$401/12800$ & $1913/2048$ & $0.000118321042$ & $0.031324476155738683$\\
$7/200$ & $1$ & $0.065210231692$ & $0.032749571$\\
$1/25$ & $3/2$ & $0.132584532412$ & $0.034762$\\
$9/200$ & $9/4$ & $0.182372427581$ & $0.036882534$\\
$1/20$ & $1277/512$ & $0.212458845389$ & $0.039482$\\
$3/50$ & $1$ & $0.236184893938$ & $0.045944547$\\
$15867/102400$ & $63/512$ & $0.294954254754$ & $0.109392080$
\end{tabular}
\end{center}
The first exact loss is $31324476155738682727829/10^{24}$.
Every successful final price is strictly below $\theta n_0$.
The set consists of a finite grid and rational local refinements;
no untested interval is excluded. For fixed positive parameters the
central bound has form $A+\beta sG_Y+M/s$ with $M\ge0$, so it is
operator convex in $s$, but its derivative has no fixed sign. It
therefore supplies no monotone bisection or global optimality result.

\section{Exact signed-action cancellation}\label{signedsupp}
The full identities in S6 and \eqref{fullbridge} give
\[
 A_Y=\mathcal E-R_q-T=SL_8S-2p_Y\otimes p_Y-T,
 \qquad T=S(J-J_8)S,\quad\|T\|\le t_8.
\]
Let $\widehat J$ be the actual smooth Fourier action from S9--S10,
and define
$\widehat Z=\widehat J-2p_Y\otimes(Y^*p_Y)$ with its exact physical
pole pairing. The continuum term has canceled before approximation.
Thus $A_YY-\widehat Z=(SL_8Y-\widehat J)-TY$ and the rational
whole-map bound is
\begin{align*}
 \|A_YY-\widehat Z\|\le{}&
 \frac{71}{100}\left(\frac{1532843}{10^{11}}+
       \frac{1797275}{10^{14}}+10^{-88}\right)+10^{-157}\\
 &+\frac{32237}{10^{12}}\frac{2159290}{10^6}<\frac{11}{10^6}.
\end{align*}
Every prime tail, Fourier tail and physical output remains.
This improves a signed-action approximation cost, not the global
form floor. The combined positive-action cache by itself does not
determine all separate smooth/pole cross Grams for this approximant.
On a larger verified quotient, both trial and approximant must be
projected. Contractivity preserves the error but not unchanged rank
or inverse claims.

\section{A cached positive five-dimensional subspace}\label{signedcache}
A different signed approximant can be formed from the existing combined
cache. Write $U=\widehat A$ and choose the supplied exact binary selector
$X$ for $N^{-1}C^*Y$. Set $Z_s=U-CX$. The full dense capture gives
$0\le R_m-CN^{-1}C^*\le(3/12500)I$, hence
\begin{align*}
 \|A_YY-Z_s\|\le{}&g+
       (3/12500+t_8)\|Y\|\\
 &+\|C\|\,\|N^{-1}C^*Y-X\|<7/12500.
\end{align*}
The selector error is paid by its exact-input solve residual,
$N\ge n_0I$, the full $CY$ uncertainty, and the whole saved-center
$N$ error at most $2\cdot10^{-8}$. Its whole physical Grams are
\[
 Y^*Z_s=YA-CY^*X,\quad C^*Z_s=CA-H_CX,
\]
\[
 Z_s^*Z_s=AG-CA^*X-X^*CA+X^*H_CX.
\]
All input and saved binary-center errors are included; in particular
$AG$ is not replaced by a projected action Gram.

Let $W_5$ be the supplied exact $56\times5$ selector and $V=YW_5$.
The enclosed physical Gram is above $0.5046238284239I$ and
$\|V\|^2<5$. The Hermitian signed center, after subtracting its
finite error and $(7/12500)\|Y\|\|W_5\|^2$, exceeds
$0.0012532198376I>I/1000$. Therefore
\begin{equation}\label{fivepos}
 V^*V>\tfrac12I,\qquad
 \mathfrak a(f)>\|f\|^2/5000\quad(0\ne f\in\Ran V).
\end{equation}
The input norm here is Euclidean until the displayed physical-norm
conversion using $\|V\|^2<5$.

Let $P_5$ be the actual projection onto $\Ran V$. A supplied exact
selector $L$ approximates the projection of $Z_sW_5$ into $V$.
Its full residual Gram is
$(Z_sW_5-VL)^*(Z_sW_5-VL)$, with every mixed term retained.
The certified lower bound for $V^*V$ justifies its inverse and prices
the difference between the selector and the exact projection.
Subtract that price and $(7/12500)\|W_5\|$ from the certified largest
residual norm for a lower bound. For an upper squared norm, Young's
inequality with parameter $1/100$ retains both the residual and
whole-map error. The exact-input residual/polar enclosures give
\begin{equation}\label{fivecoupling}
 \|(I-P_5)A_YV\|>11/50,\qquad
 \|(I-P_5)A_YV\|^2<0.062842088203.
\end{equation}
These are norms of the map with the ordinary Euclidean coefficient
input of $V$, not a normalized physical metric. A large coupling
is not a negative signed direction. Nor is $\Ran C$ proved to lie
inside this five-dimensional space, so its complement does not
inherit the $-\rho$ floor.

\section{A fully coupled Schur criterion and its scope}\label{schur}
Keep the old frame by taking $J_*= [C,Y]$ and $F=\Ran J_*$.
Every column lies in $\Dom(A_Y)$, and \eqref{rest} implies
$\mathfrak a|_{F^\perp}\ge-\rho I$ because $C^*v=0$ there.
This statement requires neither an inverse of $J_*^*J_*$ nor a
certified rank of the union.
Let $Z_*$ approximate $A_YJ_*$ and suppose
$D=A_YJ_*-Z_*$ satisfies $D^*D\le E$, with $E\ge0$.
For any exact coefficient selector $L$, put
\[
 H=J_*^*J_*,\quad B=\operatorname{Herm}(J_*^*Z_*),\quad
 Q_L=(Z_*-J_*L)^*(Z_*-J_*L).
\]
For $d>\rho$ and $\beta,\eta>0$, the sufficient comparison is
\begin{equation}\label{schurcriterion}
 B+(d-\eta/2)H-\frac E{2\eta}
 -\frac{(1+\beta)Q_L+(1+1/\beta)E}{d-\rho}\ge0.
\end{equation}
It proves $\mathfrak a(u)\ge-d\|u\|^2$ on the full score quotient.
Indeed write $u=J_*x+v$, $v\perp F$. Young's inequality gives
$\operatorname{Herm}(J_*^*D)\ge-(\eta H+E/\eta)/2$.
For the actual projection $P_F$,
$(I-P_F)Z_*=(I-P_F)(Z_*-J_*L)$, so the outside true-action Gram is
at most $(1+\beta)Q_L+(1+1/\beta)E$. Complete the square against
the complementary lower bound $(d-\rho)\|v\|^2$.
This proves \eqref{schurcriterion} even when the representation by
$J_*$ is nonunique. It is the standard Schur principle with the
required coercivity explicitly supplied, not a free unshifted inverse
\cite[Theorem 1.2]{DSS}.
For block map errors $e_C,e_Y$, one may use
$E=\operatorname{diag}((1+\lambda)e_C^2I,(1+1/\lambda)e_Y^2I)$.
The mixed signed-action Grams needed to verify this criterion for
$[C,Y]$ have not been established in the present release.

There is also a precise limitation for one declared $Y$-only capture
class. If $W\ge0$ and $YWY^*\le I$, then $YWY^*\le P_Y$.
Consequently a sufficient complement price inherited from
\eqref{rest} is
\[
 \rho_{\rm eff}=\rho+\frac{11}{10n_0}
      \|H_C-CY\,W\,CY^*\|.
\]
Let $G_c,CY_c$ be the saved centers and
$g_0=0.014115739$ exceed the accepted $G_Y$ error. The positive
matrix $M=G_c+g_0I$ is invertible, and $W_g=M^{-1}$ is the largest
capture in the cone $0\le W\le W_g$. Every smaller capture leaves
a larger positive residual. An exact binary Rayleigh locator for
the center residual gives the outward value
$0.13436493992529662392491776687409$.
With $h=32/10^{10}$, $c=1/20000$, its full perturbation is at most
\[
 h+2c\|W_gCY_c^*\|_F+c^2\|W_g\|_F<0.000384941947.
\]
The native 192-bit inverse/Rayleigh check therefore yields
$\rho_{\rm eff}>0.1478676174>0.14$ throughout that cone.
This is a bound on that sufficient certificate's price, not a lower
bound on a negative eigenvalue of the actual signed complement.

\section*{Certificate correspondence and verification boundary}
The section numbers above are the precise supplementary references in
the main article. The manifest names the immutable binary inputs,
executed-source snapshots and certificate records for each calculation.
The correspondence is as follows:
\begin{center}
\begin{tabular}{ll}
S0, S4 & symbol cells, Fourier grids, physical/zero tails\\
S1--S3 & midpoint factors, dense residual, physical norms and Gram\\
S4 & low-zero roots, Taylor grid and full sample/correlation matrices\\
S5 & lower masses and weighted deficit norm\\
S6--S7 & energy resolvent and whole matrix residual\\
S8 & equality-jet integrals and fixed-price barrier\\
S9--S10 & enriched trial, whole Fourier action and full Grams\\
S11 & rational shift endpoints and their selected matrices\\
S12--S13 & signed cancellation, cached subspace and outside coupling\\
S14 & inverse/Rayleigh robust-capture check
\end{tabular}
\end{center}
The bounded release verifier checks hashes, exact rational consequences,
selected enriched comparisons reconstructed from cached matrices, and
the stated robust-capture inverse/Rayleigh calculation. It accepts the
upstream quadrature, zero completeness, Taylor, large-Gram and whole
DFT/action allowances as certificate inputs. It does not regenerate
those expensive computations. The analytic arguments in S0--S14 explain
how those inputs control the actual continuum problem; a successful
finite replay is not by itself a formalization of those arguments.
The release verification report and formal correspondence document list
the checked Lean declarations and their explicit analytical inputs.

\begin{thebibliography}{9}
\bibitem{Theta} L. van Hemert, \emph{Sharp logarithmic coupling estimates
for theta equality frames in a signed Weil-form realization}, v1.0.0,
11 October 2026, \href{https://doi.org/10.5281/zenodo.23298499}
{doi:10.5281/zenodo.23298499}, Lemmas 2.1--2.2 and 5.1.
\bibitem{DLMF} NIST Digital Library of Mathematical Functions,
\href{https://dlmf.nist.gov/5.9.E13}{equation 5.9.13}.
\bibitem{HSW} E. Hasanalizade, Q. Shen and P.-J. Wong,
\emph{Counting zeros of the Riemann zeta function}, J. Number Theory
235 (2022), 219--241, Corollary 1.2, equation (1.5),
\href{https://www-math.nsysu.edu.tw/~pjwong/stuff/CountingRiemannZeros.pdf}
{author-hosted published article}.
\bibitem{FLINT} FLINT project, \emph{FLINT 3.6.0},
\href{https://raw.githubusercontent.com/flintlib/flint/v3.6.0/doc/source/acb_dirichlet.rst}
{versioned \texttt{acb\_dirichlet} documentation}.
\bibitem{DSS} G. Dusson, I. M. Sigal and B. Stamm,
\emph{The Feshbach--Schur map and perturbation theory},
\href{https://arxiv.org/abs/2105.02058v1}{arXiv:2105.02058v1}, Theorem 1.2.
\end{thebibliography}
\end{document}
